\noindent\hspace*{1.2em}\kbvlabel{ಉದಾ}\hspace{0.6em}1. ಬೆಕ್ಕು ಇಲಿಯನ್ನು ಕೊಂದಿತು.\par

\indent {[}ಇಲಿ + ಅನ್ನು = ಇಲಿಯನ್ನು ; ಇದು ಯಕಾರಾಗಮಸಂಧಿ{]}\par

\vspace{2.5pt plus 3pt minus 0.5pt}\noindent\hspace*{1.2em}\kbvnum{2.}\hspace{0.55em}ಕೋತಿಯು ಹಣ್ಣನ್ನು ತಿಂದಿತು.\par

\vspace{4pt plus 4pt minus 1pt}\noindent{\kbvsecfont\fontsize{11pt}{14.5pt}\selectfont\color{kbvaccentdark} (3)\hspace{0.45em}ತೃತೀಯಾ ವಿಭಕ್ತಿ :}\hspace{0.5em}ತೃತೀಯಾ ವಿಭಕ್ತಿ ಪ್ರತ್ಯಯ “ಇಂದ” ಆಗಿದೆ.\par

\noindent\hspace*{1.2em}\kbvlabel{ಉದಾ}\hspace{0.6em}1. ಸಕ್ಕರೆಯಿಂದ ಮಿಠಾಯಿ ತಯಾರಿಸುತ್ತಾರೆ.\par

\vspace{2.5pt plus 3pt minus 0.5pt}\noindent\hspace*{1.2em}\kbvnum{2.}\hspace{0.55em}ಹಣದಿಂದ ಗುಣ ಸಂಪಾದಿಸಲು ಸಾಧ್ಯವಿಲ್ಲ.\par

\vspace{2.5pt plus 3pt minus 0.5pt}\noindent\hspace*{1.2em}\kbvnum{3.}\hspace{0.55em}ಅಭ್ಯಾಸ ಮಾಡಿ ಮುಂದುವರೆದಿದೆ.\par

\vspace{2.5pt plus 3pt minus 0.5pt}\noindent\hspace*{1.2em}\kbvnum{4.}\hspace{0.55em}ಮೂರಕ್ಕೆ ಐದು ಸೇರಿದರೆ ಎಂಟಾಗುವುದು.\par

\vspace{4pt plus 4pt minus 1pt}\noindent{\kbvsecfont\fontsize{11pt}{14.5pt}\selectfont\color{kbvaccentdark} (4)\hspace{0.45em}ಚತುರ್ಥಿ ವಿಭಕ್ತಿ :}\hspace{0.5em}ಚತುರ್ಥಿ ವಿಭಕ್ತಿ ಪ್ರತ್ಯಯಗಳು “ಇಗೆ, ಗೆ, ಕ್ಕೆ, ಅಕ್ಕೆ” ಆಗಿವೆ.\par

\noindent\hspace*{1.2em}\kbvlabel{ಉದಾ}\hspace{0.6em}1. ರಾಮನಿಗೆ ಕಾಲು ಉಳುಕಿದೆ.\par

\vspace{2.5pt plus 3pt minus 0.5pt}\noindent\hspace*{1.2em}\kbvnum{2.}\hspace{0.55em}ಆನೆಗೆ ನಾಲ್ಕು ಕಾಲುಗಳುಂಟು.\par

\vspace{4pt plus 4pt minus 1pt}\noindent{\kbvsecfont\fontsize{11pt}{14.5pt}\selectfont\color{kbvaccentdark} (5)\hspace{0.45em}ಪಂಚಮೀ ವಿಭಕ್ತಿ :}\hspace{0.5em}ಪಂಚಮೀ ವಿಭಕ್ತಿ ಪ್ರತ್ಯಯ “ದೆಸೆಯಿಂದ” ಎಂಬುದಾಗಿದೆ.\par

\noindent\hspace*{1.2em}\kbvlabel{ಉದಾ}\hspace{0.6em}1. ಹಣದ ದೆಸೆಯಿಂದ ಸ್ನೇಹ ಮುರಿದು ಬಿತ್ತು.\par

\vspace{2.5pt plus 3pt minus 0.5pt}\noindent\hspace*{1.2em}\kbvnum{2.}\hspace{0.55em}ಕೌರವರ ದೆಸೆಯಿಂದ ಯುದ್ಧ ಸಂಭವಿಸಿತು.\par

\vspace{4pt plus 4pt minus 1pt}\noindent{\kbvsecfont\fontsize{11pt}{14.5pt}\selectfont\color{kbvaccentdark} (6)\hspace{0.45em}ಷಷ್ಠಿ ವಿಭಕ್ತಿ :}\hspace{0.5em}ಷಷ್ಠಿ ವಿಭಕ್ತಿ ಪ್ರತ್ಯಯ “ಅ” ಆಗಿದೆ.\par

\noindent\hspace*{1.2em}\kbvlabel{ಉದಾ}\hspace{0.6em}1. ರಾಜನ ಹೆಂಡತಿ ರಾಣಿ.\par

\vspace{2.5pt plus 3pt minus 0.5pt}\noindent\hspace*{1.2em}\kbvnum{2.}\hspace{0.55em}ದುರ್ಜನರ ಸಂಗ ಮಾಡಬಾರದು.\par

\vspace{4pt plus 4pt minus 1pt}\noindent{\kbvsecfont\fontsize{11pt}{14.5pt}\selectfont\color{kbvaccentdark} (7)\hspace{0.45em}ಸಪ್ತಮೀ ವಿಭಕ್ತಿ :}\hspace{0.5em}ಸಪ್ತಮೀ ವಿಭಕ್ತಿ ಪ್ರತ್ಯಯ “ಅಲ್ಲಿ” ಆಗಿದೆ.\par

\noindent\hspace*{1.2em}\kbvlabel{ಉದಾ}\hspace{0.6em}1. ನನ್ನ ಜೇಬಿನಲ್ಲಿ ಹಣವಿಲ್ಲ.\par

\vspace{2.5pt plus 3pt minus 0.5pt}\noindent\hspace*{1.2em}\kbvnum{2.}\hspace{0.55em}ದೇವರಲ್ಲಿ ಭಕ್ತಿಯಿಡು.\par
